% ============================================================================
% Section 4: Formalization
% ============================================================================

\section{Formalization}\label{sec:formalization}

Before presenting our contributions in detail, we formalize the notions of \textbf{Embodied Agent Environments (EAEs)} and \textbf{Agentic Harnesses}, which provide the vocabulary for describing our benchmark (Chapter~\ref{sec:benchmark}), harness architecture (Chapter~\ref{sec:pokeagent}), and AutoEvolve algorithm (Chapter~\ref{sec:autoevolve}). We then define a \emph{harness spectrum} that frames the central experimental question of this thesis.

\subsection{Embodied Agent Environments}\label{sec:formal_env}

We consider an embodied agent interacting with a visual environment through a minimal interface. At each discrete timestep $t$, the agent receives a frame observation $o_t$ (a rendered image of the current game state) and a structured text state $x_t$ (derived from read-only memory and map data), then selects an action $a_t$ from a fixed set of button inputs $\mathcal{A}$. The environment is partially observable: the agent cannot access the full internal game state encompassing a global view of the entire environment. In the POMDP formulation introduced in Section~\ref{sec:mdp}, the pair $(o_t, x_t)$ constitutes the agent's observation, and the hidden state $s_t$ resides in GBA RAM.

Formally, we model the environment as a partially observable process:
\begin{equation}\label{eq:env}
    (o_{t+1}, x_{t+1}) = \mathcal{E}(s_t, a_t), \quad s_{t+1} = f(s_t, a_t),
\end{equation}
where $s_t \in \mathcal{S}$ is the hidden game state, $f$ is the deterministic (up to random encounters) state transition function implemented by the game engine, and $\mathcal{E}$ is the observation function that renders the visible portion of the state. This formulation is general: any environment that exposes visual output and accepts discrete inputs fits this interface, making our framework applicable beyond Pok\'{e}mon to other visual, embodied domains.

The agent maintains a trajectory $\tau_t = \{(o_1, x_1, a_1), \ldots, (o_t, x_t, a_t)\}$ recording its full interaction history. In practice, this trajectory also includes the agent's reasoning traces (chain-of-thought outputs) and tool-call records at each step, forming a rich log that serves as the substrate for context management, reflection, and evolutionary optimization.

\subsection{Agentic Harnesses}\label{sec:formal_harness}

An \emph{agentic harness} $\mathcal{H}$ is the scaffolding layer between a foundation model $M$ and the environment. Following the decomposition introduced in the \pokeagent{} Challenge~\citep{karten2025pokeagent}, a harness mediates agent behavior through five components:
\begin{itemize}[leftmargin=*]
    \item \textbf{System prompt} $p$: instructions and strategic guidance provided to the model at each reasoning step.
    \item \textbf{Subagents} $\mathcal{S} = \{s_1, \ldots, s_k\}$: specialized modules that can be invoked by the orchestrator for specific tasks (e.g., battle strategy, navigation planning, self-reflection).
    \item \textbf{Memory} $\mathcal{M}$: a persistent knowledge store that accumulates facts, strategies, and observations across the agent's trajectory.
    \item \textbf{Skills} $\mathcal{K}$: reusable behavioral routines, either textual (strategic guidance) or executable (Python code in a sandbox), that encode solutions to previously encountered situations.
    \item \textbf{Tools} $\mathcal{T}$: programmatic capabilities available to the model, such as pathfinding, information lookup, or code execution.
\end{itemize}

A harness configuration is thus a tuple $\mathcal{H} = (p, \mathcal{S}, \mathcal{M}, \mathcal{K}, \mathcal{T})$, and the agent's behavior at each step is determined by $a_t = M(o_t, x_t, \tau_{<t}; \mathcal{H})$, where $M$ denotes the foundation model's inference conditioned on the current observation, history, and harness state.

\subsection{Harness Spectrum}\label{sec:harness_spectrum}

A central claim of this thesis is that harness design is a first-order variable in agent performance, often mattering as much as the choice of underlying model. To make this claim precise, we define a spectrum of harness configurations over the five-component tuple. Because several harness components accumulate state during play (the agent creates memories, learns skills, and may spawn subagents at runtime), we distinguish between the \emph{initial} configuration at session start (subscript~$0$) and the \emph{evolved} configuration after extended play (subscript~$1$). Importantly, $\mathcal{M}_0$ and $\mathcal{K}_0$ are empty for all harnesses; the difference between configurations lies in the system prompt, the pre-defined subagents and tools, and how these components grow during execution.

\begin{itemize}[leftmargin=*]
    \item \textbf{Minimal harness} $\mathcal{H}_{\min}$. Initialized as $(p_{\min,0},\, \mathcal{S}_{\min,0},\, \mathcal{M}_0,\, \mathcal{K}_0,\, \mathcal{T}_{\min})$, where $p_{\min,0}$ is a default system prompt, $\mathcal{S}_{\min,0}$ exposes subagent primitives (the ability to create and invoke subagents) but includes no pre-built custom subagents, $\mathcal{M}_0 = \mathcal{K}_0 = \varnothing$, and $\mathcal{T}_{\min}$ contains only the base tools required to interact with the environment (send input, capture frame, read state) and the persistent stores. During play, the agent may accumulate state:
    \[
        \mathcal{H}_{\min} \;\longrightarrow\; (p_{\min},\, \mathcal{S}_{\min,1},\, \mathcal{M}_1,\, \mathcal{K}_1,\, \mathcal{T}_{\min}).
    \]

    \item \textbf{Expert harness} $\mathcal{H}_{\text{expert}}$. Initialized as $(p^*_0,\, \mathcal{S}^*_0,\, \mathcal{M}_0,\, \mathcal{K}_0,\, \mathcal{T}^*)$, where $p^*_0$ is a refined system prompt encoding domain strategy, $\mathcal{S}^*_0$ includes purpose-built subagents (battle, reflection, verification, planning), and $\mathcal{T}^*$ is the extended tool suite including A* pathfinding, research tools, and code execution. This is the \pokeagent{} system described in Chapter~\ref{sec:pokeagent}. During play:
    \[
        \mathcal{H}_{\text{expert}} \;\longrightarrow\; (p^*_0,\, \mathcal{S}^*_1,\, \mathcal{M}_1,\, \mathcal{K}_1,\, \mathcal{T}^*).
    \]

    \item \textbf{Auto-evolved harness} $\mathcal{H}_{\text{auto}}$. Initialized identically to $\mathcal{H}_{\min}$: $(p_{\min,0},\, \mathcal{S}_{\min,0},\, \mathcal{M}_0,\, \mathcal{K}_0,\, \mathcal{T}_{\min})$. AutoEvolve (Chapter~\ref{sec:autoevolve}) iteratively modifies all five components through in-context evolutionary optimization over trajectory segments, without environment resets:
    \[
        \mathcal{H}_{\text{auto}} \;\longrightarrow\; (p_{\text{auto},1},\, \mathcal{S}_{\text{auto},1},\, \mathcal{M}_{\text{auto},1},\, \mathcal{K}_{\text{auto},1},\, \mathcal{T}_{\min,1}).
    \]
\end{itemize}

An important subtlety is that $\mathcal{H}_{\min}$ and $\mathcal{H}_{\text{expert}}$ share the same primitive access to subagents, memory, and skills; their difference lies in \emph{what is pre-built} (custom subagents, curated prompts, extended tools) rather than in the presence or absence of these capabilities. Both harnesses can accumulate $\mathcal{M}_1$ and $\mathcal{K}_1$ through play. This overlap means that performance differences between $\mathcal{H}_{\min}$ and $\mathcal{H}_{\text{expert}}$ reflect the value of human-engineered initialization rather than a strict ablation of harness components. We revisit the implications of this design choice in the discussion of Chapter~\ref{sec:autoevolve} and identify the development of a more isolated minimal baseline as a direction for future work (Chapter~\ref{sec:outlook}).

The central experimental question is whether $\mathcal{H}_{\text{auto}}$, starting from $\mathcal{H}_{\min}$, can approach the performance of $\mathcal{H}_{\text{expert}}$. Answering this requires evaluating which harness components are most amenable to automated discovery, how quickly the evolutionary process converges, and whether auto-evolved configurations learn strategies and abstractions that transfer across environment resets. Chapter~\ref{sec:autoevolve} presents the AutoEvolve algorithm and evaluates these questions empirically.
